\chapter[Life without Ready-Made Meaning]{What If Life Has No Ready-Made Meaning?}
\section{An Unsent Postcard}
Begin with something small: a postcard.

It fell from one of Grandad's old books on moving day. The brown-paper cover had grown brittle; the card lay around page two hundred, its corners worn round. The photograph had faded to pale yellow-green: a lakeside park decades ago, with a coin-operated telescope of a kind long since disappeared.

Turn it over. A stamp is attached; the recipient's address is stamped in the upper-right corner. The handwriting is neat with a young person's excessive care. The message runs only two and a half lines:

\begin{quote}
Old Zhou, greetings. The lakeside has changed a lot, but the telescope is still here. When your summer holiday begins...
\end{quote}
The sentence stops there. Nothing follows the holiday, no signature, no postmark. It was never sent.

Holding it, you wonder: who was Old Zhou? Did his summer holiday come? How would the sentence have ended: ``let's go again'' or ``there are things I want to say in person''? Why did the writer stop? Was he called to dinner, did he decide there was no need to write, or did something happen that permanently removed the chance to finish?

A stranger thought follows: do those two and a half lines count? They reached nobody and changed nothing. Apart from you holding the card now, nobody knows it exists. Were the ten minutes spent writing it a real part of that person's life? If a whole life ends like an unsent postcard, with unfinished deeds and undelivered words, do its sentences still count?

Keep the postcard beside you. This chapter's questions rest beneath those two and a half lines. Life comes with no factory-printed meaning, and everything may ultimately remain ``unsent.'' Why, then, do we write?

\section{Laughter at a Funeral}
Now come into a scene.

Time: nine on a Saturday morning in early November. Place: a funeral hall in a small southern city. Fourteen-year-old A He stands among the mourners, a black gauze armband scratching her skin. Three smells mingle: incense smoke, chrysanthemums, and damp rising from a heavily trodden carpet.

At the front stands her grandmother's portrait, taken last year: a broad smile showing uneven teeth. The celebrant reads the eulogy in an excessively level voice. At ``a lifetime of hard work,'' A He hears her mother sob on her left. Her aunt passes tissues; the packet's tearing sounds harsh in the quiet hall.

Afterwards, relatives eat beneath a courtyard canopy, where six or seven round tables stand in bright November sun. Then something makes A He's heart jolt: someone tells a joke and a whole table laughs, loudly. Her great-uncle booms out a finger-guessing drinking game at the next table. Younger cousins finish eating and run off, chasing one another through white parking-space lines, laughing louder than the adults.

A He stands at the canopy's edge holding her bowl, anger rising from her stomach. Grandma is dead. Dead. How can you laugh?

During the evening vigil she tells her uncle. He is adding spirit money to a brazier, firelight on his face. After thinking, he says, ``Your grandma hated a quiet house. At New Year meals, if anyone put down their chopsticks and left early, she complained for ages. With all this liveliness today, she would be the happiest person here.''

A He is not entirely convinced but stops arguing. Later, relatives doze against chair backs and incense burns down to its base. Staring at the portrait, she thinks of something colder than their laughter:

Next week Grandma's room will be cleared. Her cup, reading glasses, and pickling jars will be dealt with one by one. In a few years, beyond occasional family mentions, nobody will remember her dislike of quiet or her very salty pickled radishes. Decades later, even those who remember will be gone. Then ``Grandma hated a quiet house'' will hang without a recipient, like the postcard's unfinished sentence.

What about me? A He pulls her coat tight. Will my exams, memorised vocabulary, and collected stickers end the same way?

\section{The Question That Appears at Night}
State the conflict before rushing to an answer.

On Monday after the funeral, A He's life fits perfectly back into its old track: morning reading, group exercise, the weekly maths test. Nothing in the world changes its schedule for Grandma's absence. On Tuesday evening, her pen stops halfway through English vocabulary. A question catches up from the late-night vigil and sits squarely on her exercise book:

\textbf{If everything eventually disappears, what is the point of doing this now?}

This is neither adolescent posturing nor disguised exam anxiety. It is among humanity's oldest cracks. For millennia, people from palaces to huts, philosophers and shepherds alike, have stopped before it. Sooner or later you will encounter it at an ordinary moment: after a failed test, during a sleepless night, or when everything seems strangely smooth.

Two common, inadequate reactions surround this crack. Identify them first.

\textbf{First reaction: mistake the question for an answer.} ``Exactly. Everything disappears, so nothing is worthwhile.'' This sounds profound but is counterfeit. It does not answer the question; it borrows its force to announce that the show is over. Later we will see that this answer does not even believe itself.

\textbf{Second reaction: treat the question as a mistake.} ``Thinking about that is useless. Study properly.'' Adults often say this, partly from kindness: they fear the crack will swallow you. But ask what ``useless'' means. Whose ledger determines usefulness? If the ledger itself, the standard defining use, is precisely what the question challenges, calling the question useless is not answering but refusing to enter the discussion.

Only now does the real problem take shape, in two uncomfortably hard parts:

First, life includes no instruction manual. At birth nothing tells you what you came to do. The universe does not supply purposes; science cannot help here, and prayer cannot seal off doubt.

Second, life has a deadline. However much you write, it may stop at ``when your summer holiday begins.''

Together they form our question: \textbf{if life has no ready-made meaning, what should we do?}

Do not rush ahead for answers. Stand beside A He's desk first. If everything disappears, where can reasons for living well now come from?

\section{Four Answers on the Stage}
How old is this question? Old enough for humanity to develop several complete sets of answers. Let each take the stage and finish speaking, including those you may dislike.

\textbf{Voice one: meaning is given, and that is good.}

This has been humanity's predominant answer. In Christian tradition, meaning comes from God's plan: believers understand themselves as created, loved, and awaited. In Confucian tradition, meaning belongs to a long chain: honour ancestors and raise descendants, receiving one baton and passing the next. In Buddhist tradition, permanence and self are illusions; seeing through them and escaping attachment offers release.

State this position fully. Dismissing it with ``What century is this?'' is unfair.

Meaning supplied by something larger solves at least three real problems. First, \textbf{suffering has somewhere to belong}. People most fear not pain but pointless pain. Belief that things form part of a plan places hospital or battlefield suffering inside a larger story one can bear. ``Invent your own meaning'' often sounds weightless before real disaster. Second, \textbf{you need not reinvent the wheel alone}. Requiring every ordinary person to prove life's worth from scratch is as cruel as requiring everyone to invent mathematics. Tradition hands over answers tested and refined by hundreds of millions: this is civilisation's basic operation, not laziness. Third, \textbf{a community comes with it}. Meaning cannot be sustained alone in a room; it needs festivals, rituals, meals, and relatives sharing a vigil. Given meaning comes with fellow believers.

This is a strong case. Its power lies within the funeral laughter. Those relatives endured the weekend not through philosophical proofs but through inherited practices: adding spirit money, setting tables, keeping vigil. These spoke words they could not formulate: she has gone; we remain.

\textbf{Voice two: nothing is given, so there is nothing.}

Nihilism answers that the universe does not care. Stars ignore your monthly exam; the galaxy knows nothing of Grandma's dislike of solitude; centuries hence nobody will remember you. Stop pretending: nothing is worthwhile.

Its case rests on truthfulness. It offers no comforting fairy tale and pretends to see no script. That honesty deserves respect. But it contains a fatal flaw we will examine later: compressing three different questions into one sentence.

\textbf{Voice three: the sense of meaning is a brain mechanism.}

Biology offers this description: brains that experience meaning are better at surviving and raising offspring, so natural selection installed a meaning generator. Finding sunsets magnificent or babies worth caring for essentially comes off the same production line as wanting food when hungry.

This may be correct, but distinguish its level. It explains \textbf{where a sense of meaning comes from}, a factual question, not \textbf{whether meaning counts}. Evolutionary hunger does not make eating an illusion. Throughout this chapter, distinguish what happened, why, how people understood it, and how we evaluate it. None can impersonate another.

\textbf{Voice four: no ready-made answer means freedom begins.}

This voice arrived latest historically but stands closest to your life. Its speakers are called \textbf{existentialists}: provisionally, philosophers insisting that people must answer for themselves. Meet four representatives:

\textbf{Søren Kierkegaard}, a nineteenth-century Dane, focused on \textbf{anxiety}. Broadly, fear has an object, such as dogs or exams, while anxiety does not. The dizziness at a cliff's edge comes not from ``I may fall'' but from ``I could jump'': the first recognition that \textbf{you have a choice}. His surprising conclusion makes anxiety not a malfunction but freedom's admission ticket. More possibilities mean greater dizziness. Someone never anxious may simply not realise they stand at a cliff.

\textbf{Friedrich Nietzsche}, a late-nineteenth-century German, declared ``God is dead'' as an alarm, not a celebration. Europe's thousand-year pillar of meaning was collapsing, and most people had not grasped the consequences. Old answers were gone; new ones would not drop from the sky. Humanity must collapse into emptiness or learn to \textbf{create values}. ``Become who you are'' meant not discovering a finished inner self but sculpting yourself, blow by blow, from the available material.

\textbf{Jean-Paul Sartre}, a twentieth-century Frenchman, declared that \textbf{existence precedes essence}. A paper knife is designed, its essence established, before it is made and exists: its purpose precedes it. Humans reverse this: thrown into the world first, we determine what we are through choices. No manual of human nature guarantees us. Hence his heavier claim: we are \textbf{condemned} to freedom. Freedom is a sentence, not a benefit, with nowhere to appeal. ``I had no choice; everyone does it'' fails: following the crowd was itself your choice.

\textbf{Albert Camus}, a Frenchman born in Algeria, supplies our first keyword, \textbf{the absurd}. This means not comedy but the confrontation between humanity's cry for meaning and a silent world. It resides neither in humanity nor in the world alone, but in their gap. Camus takes Sisyphus as his symbol: punished to push a huge stone uphill, watch it roll back, descend, and repeat forever. Like your vocabulary lists and school exercise routines, another always follows. Yet Camus reaches a surprising conclusion: \textbf{we must imagine Sisyphus happy}. Descending with full knowledge that the stone will fall again, yet continuing, he makes the stone and his fate his own rather than the gods'.

Correct an easy misconception. People often call existentialists gloomy thinkers who permit anything you feel like doing. \textbf{This is an error serious enough to require explicit correction.} Sartre's freedom is a burden, not a licence. ``Whatever I feel like'' evades responsibility by blaming desire. Camus opposed both ending one's life as surrender to absurdity and false hope pretending absurdity absent. He sought a third position: remaining lucidly within it without leaving. Existentialism is not an elegant version of gloom culture; among these answers, it insists most fiercely on responsibility.

\section[Thinking Bridge: Three Questions]{Thinking Bridge: Split Meaninglessness into Three Questions}
Voice two left a promise: nihilism has a fatal flaw. Now examine it with a portable tool.

The reminder is this: \textbf{``Life has no meaning'' flattens three separate questions into one.} Whenever you hear or think it, ask which level has collapsed.

\textbf{Level one: does the universe have a purpose?}

Is the universe going somewhere? Does a master script assign everything a use? The honest answer is that we do not know, and very likely not. Science describes how the universe operates, not what it is for. If this level collapses, leave it: you cannot defend it, and it is not your responsibility.

\textbf{Level two: is my life worth living?}

Are my few decades, overall, something worth having happened? This level \textbf{does not follow from the first}. The absence of a cosmic script does not make your match worthless. Football is absent from that script too, yet players can be moved to tears. Scripts are one thing, matches another.

\textbf{Level three: is there something worth doing now?}

Is anything concrete worth my effort today, at this moment? Its connection to the other levels is looser still. Someone may firmly believe in cosmic purpose, with level one intact, yet find nothing interesting after the university entrance examination, a level-three collapse. Conversely, someone convinced of cosmic silence may enthusiastically divide orchids or mark an apprentice's work every day: level one collapsed, level three working well.

The tool immediately reveals A He's night-time question as a \textbf{misdiagnosis}. What collapsed was level three: after Grandma's death, memorising words suddenly seemed insignificant. In the dark she mistook this for level one, as if permission to continue depended on telephoning the universe for its overall purpose. That call never connects, leaving her frozen. Most outbreaks of life's meaninglessness are level-three faults reported as level-one disasters.

How do we repair level three? Human experience offers four building materials, four sources of meaning:

\begin{itemize}
\item \textbf{Discovery.} Meaning is sometimes found, like a road already waiting, rather than invented. In Nazi concentration camps, psychologist Viktor Frankl observed that those lasting longest often held onto something awaiting completion outside: an unfinished manuscript, a person they had not seen. They did not fabricate it but recognised it as already theirs.
\item \textbf{Creation.} Choose something and invest effort; meaning grows with that investment. After three years of practice, the piano becomes part of your bones.
\item \textbf{Inheritance.} Grandma's radish recipe, New Year customs, your language: meanings made by others and passed to you. Accept them and they become yours.
\item \textbf{Shared maintenance.} A team's traditions, a class's annual autumn outing, a family's unfailing New Year vigil belong to nobody's sole authority. Everyone renews the subscription together each year. The funeral meal was one such renewal.
\end{itemize}
Fourth-level collapse belongs to the universe, beyond your control; third-level collapse belongs to you and can be repaired. That is the whole tool.

\section{The Question Confucius Did Not Answer}
Read a short primary source. The question appearing at night was asked face to face in a Chinese classroom over two millennia ago.

The Analects chapter ``Xian Jin'' records Zilu, also called Jilu, asking Confucius how to serve spirits. Confucius replies:

\begin{quote}
The Master said, ``Not yet able to serve people, how can you serve spirits?''
\end{quote}
Zilu ventures further: ``May I ask about death?'' Confucius says:

\begin{quote}
``Not yet understanding life, how can you understand death?''
\end{quote}
In other words: how can death be understood before the business of living?

At first, many readers think Confucius is dodging his student's deepest question. Look closely at the redirection. He says neither that death is unreal, offering comfort, nor that an afterlife awaits, making a promise, nor ``do not ask,'' silencing inquiry. He establishes an order: life comes first and is something you can work on, serving living people and living today properly. Leave death to time and life to yourself. In our terms, Confucius returns the questioner steadily from level one to level three.

Centuries later, another Chinese thinker goes further. In Zhuangzi's ``Perfect Happiness,'' Zhuangzi's wife dies. His friend Hui Shi arrives to offer condolences and finds him sitting with legs spread, drumming on an earthenware basin and singing. Hui Shi protests furiously: she shared your life and raised children; not crying is bad enough, but singing is outrageous. Zhuangzi replies, broadly: of course I grieved at first. But originally she had neither life, form, nor breath. Through transformations amid the indistinct came breath, form, and life; now another transformation returns her, like the four seasons. She sleeps peacefully in heaven and earth's great house. Wailing would show my failure to understand destiny. So I stopped crying.

Later still, Eastern Jin poet Tao Yuanming wrote his own elegies in advance. The third of his ``Imitations of Funeral Songs'' ends:

\begin{quote}
What remains to say when dead? Entrust the body to the hills.
\end{quote}
That is: after death, what more can be said? Simply give the body to the hills.

Together these passages share a stance: \textbf{these Chinese traditional answers do not promise heaven to dispel fear; they fold death back into this life and this world}. Confucius folds it into serving people, Zhuangzi into the seasons, Tao Yuanming into the hills. Confucianism also offers a more active route. The Zuo Commentary, Duke Xiang's twenty-fourth year, says:

\begin{quote}
Highest is establishing virtue; next, establishing achievements; next, establishing words.
\end{quote}
Virtue, achievements, and words speak after their maker is gone. These are the ``three immortalities'': immortality certified not by heaven but by what remains effective in this world.

Following our rule, do not reduce the tradition to one voice. Confucius centres human affairs; Zhuangzi thinks you grip them too tightly. Confucians pursue lasting reputation through three immortalities; Daoists would have you relinquish reputation too. Their two-thousand-year argument shows how serious the crack is: nobody has completely filled it.

\section{One Wasteland, Three Ways Forward}
Lay all the parts on the table. The facts are shared: people die, the universe supplies no manual, and Grandma's room will be cleared next week. Our four distinctions remain: what happened, a fact; why and where to go, contested explanations; how participants understood it, as in the uncle's remark; and our evaluation, such as whether funeral laughter is right. Here we compare the second category: three different, nonequivalent responses to the same facts.

\textbf{Explanation one: disappearance means zero; honesty admits it.}

This is nihilism's fullest version. Its strength is honesty: no invented heaven, reserved cosmic seat, or reassuring future used as pain relief. Our tool exposed its limitation; now sharpen the point: it \textbf{cannot stand by its own behaviour}. Someone declaring that everything amounts to zero still rises tomorrow, dodges an oncoming electric scooter, and brings medicine to a feverish friend. If zero were truly everything, why dodge? Their mouth inhabits level one, cosmic purposelessness; their body inhabits level three, something worth protecting now. The position is less wrong than uninhabitable: a house built only for visitors.

\textbf{Explanation two: meaning must be created personally, with unavoidable responsibility.}

Existentialism offers the only route, it argues, that neither lies about a manual nor declares everything over. No script? Write it, and answer for every word. Its real limitation is \textbf{its weight}. All choices and responsibility being yours may liberate Sartre's reader but oppress a fourteen-year-old unable even to choose dinner. Is every mistake entirely mine? Frankl offers a correction: meaning is not always \textbf{invented} from nothing; often it is \textbf{discovered}, already waiting for recognition. Pure creation imagines lone heroes; most actual people are not.

\textbf{Explanation three: meaning already operates in relationships; continue it rather than invent it.}

This view sees the funeral differently. Laughter is neither indifference nor performance but \textbf{maintenance}. Adding spirit money, setting tables, keeping vigil, and telling the deceased's jokes jointly enact: she has gone; we remain; she hated a quiet house, and we remember. Nobody need invent meaning. It has operated for decades in the pickle recipe and the rule against leaving New Year dinner early; your task is to receive it. This closely fits ordinary life: most people survive bereavement through these practices, not philosophy. But its limitation matters equally: \textbf{inherited meaning may be someone else's script}. Honouring ancestors has crushed many lives it prescribed. Some traditions should be discarded. Pure inheritance cannot answer what to do when what is handed down is wrong.

Each route holds some truth and misses the whole. Comparing them is not declaring all correct. Nihilism fails to believe itself; creation imagines people too solitary; inheritance imagines tradition too innocent. A workable answer probably borrows responsibility from existentialism, continuity from tradition, and discovery from Frankl. That combination is no standard answer, but at least it offers an inhabitable house.

\section[Further Challenge: Facing Death]{Further Challenge: Face to Face with Death}
Now bring in the heaviest stone. Absurdity, anxiety, and choice ultimately lead to \textbf{death}, discussed by philosophers for over two thousand years. Let two major approaches seriously contend.

\textbf{First: Epicurus's argument---death is nothing to you.}

The ancient Greek Epicurus argues, broadly: death is nothing to us because all good and bad exist in sensation, while death ends sensation. When we exist, death is absent; when death exists, we are absent. We never meet it, so why fear it?

Like a polished glass sphere, the argument is smooth and tightly joined, comforting people for millennia. Squeeze it and a crack appears: it addresses fear of being dead, not fear of losing life. A He feared not Grandma's posthumous discomfort but the disappearance of memories of her. The postcard saddens us not because of what happened after its writer died, but because its sentence \textbf{never arrived}. Often we fear loss rather than the state of death: not seeing a sister grow up, the next episode, or speaking words before time runs out. Epicurus's sphere cannot roll into this room of loss. Remember: \textbf{however rigorous an argument, if it misses where people actually hurt, its rigour is merely tidiness.}

\textbf{Second: Heidegger's reminder---being-towards-death.}

Twentieth-century German philosopher Martin Heidegger reverses the approach: use mortality rather than prove death harmless. Ordinarily we live by ``everyone does it'': schooling because others attend, anxiety about house-buying because others buy. Such lives belong to the crowd rather than ourselves. Death uniquely \textbf{cannot be done for you by anyone else}. Others can do your homework or apologise for you, not die your death. Truly recognising this cracks the shell of conformity: this life is mine to live, with a deadline that prevents endless postponement. Roughly, being-towards-death makes death not a full stop but the pen compelling you to make life yours. The deadline is not a bug; it makes writing serious. The postcard's possible failure to arrive makes its two and a half lines precious.

But listen carefully to this chapter's most important qualification: \textbf{all this is philosophical language, and actual grief does not respond to it.}

Imagine your best friend, whose father has died, sitting beside you with red eyes. Reciting Epicurus---``when we exist, death is absent''---is not merely unhelpful but cruel. The pain is not a concept: Dad will not collect your friend tomorrow. Arguments can analyse conceptual death; actual grief does not enter their contest.

Grief needs another language: \textbf{companionship, ritual, and repeated remembrance}. Spirit money added at a vigil, ``Grandma hated a quiet house,'' a customary dish on a death anniversary. These do not explain death but transform something unbearable into something people can pass through together. Consider Mexico's Day of the Dead in early November: households lay marigold paths, offer the deceased's favourite tamales and drinks, play music in cemeteries, and talk all night. Tradition holds that souls return to visit. Rather than debate death, they \textbf{maintain relationships with the dead}. This speaks the same language as the funeral meal and Confucius's return to serving people.

This section leaves a distinction, not an answer: \textbf{philosophising about death and accompanying someone in grief are different crafts requiring different languages}. Answering tears with arguments uses the wrong tool; evading thought with tears also gives something up. Maturity means knowing which craft the moment requires. More practically, if you or someone nearby remains weighed down by meaninglessness, unable to eat or sleep and losing interest in everything, do not carry that with philosophy: seek a school counsellor or doctor. Existential emptiness and depression requiring support also need different responses.

\section{When Gloom Becomes Fashionable}
This ancient question now appears in your phone's newest packaging.

Social platforms offer ready-made tones: ``the world isn't worth it,'' ``I'm emo,'' ``let it slide,'' ``lying flat,'' and ``Buddha-like detachment.'' Gloom becomes slang and social currency, with the first person to seem indifferent gaining an edge. Our analytical tool separates genuine exhaustion from school or family, existential questions, and simple imitation. Calling all of it youthful profundity misreads it; calling all of it melodrama does too. Real cracks admit cold air here.

Consider a less visible mechanism: \textbf{algorithms are today's largest wholesalers of meaning}. Videos and games never ask whether life has meaning; they supply level three directly, a reason for this moment: the next clip, round, or notification. But wholesale reasons do not sustain you. Greater emptiness after three hours of scrolling is level three giving up its lease. Entertainment is not sinful, but moments supplied by others cannot build your own house.

School offers another typical collapse. From primary school through the final secondary year, many goals are \textbf{outsourced}. Parents and teachers prefabricate meaning: a good middle school, high school, university. Clear and useful, this goal has one flaw: it is not yours. Trouble comes either upon success, ``what next?'', or failure, ``what am I worth now?'' Again a level-three prefabricated goal expires but feels like cosmic collapse because its holder never practised creating level three. The university phenomenon called ``empty-heart syndrome,'' lacking nothing yet finding nothing meaningful, mostly arises this way.

New variables arrive too. Machines write essays, paint, and solve problems, leaving a new question at the door: if machines can do everything, what is left for me? You already know how to separate its levels. Machine problem-solving does not make your problem-solving worthless, just as calculators did not stop maths contestants being moved to tears.

Existentialism's practical advice today sounds disappointingly plain: \textbf{do not find the answer; give your answer}. You need not solve the universe in your head before receiving permission to act. Daily choices---finishing memorisation, defending a classmate mocked in a chat, learning Grandma's pickle recipe---vote for the person you are becoming. Meaning is not a buried coin awaiting excavation. It is voted into being, discovered, inherited, and renewed together.

\section[Your Question: An Answer Machine?]{Your Question: Would You Use an Answer Machine?}
Return to the postcard, but first consider a machine.

Imagine scientists fifty years hence creating a ``meaning machine.'' Wear its helmet for three minutes; it scans your brain and prints a personalised Life's Meaning Manual. Its goals fit your nature perfectly. Follow them and lifelong certainty and fulfilment are guaranteed; the night-time question never knocks again. Money back if it fails.

Would you wear it?

Weigh both sides fully before answering.

Reasons to accept are powerful: freedom from anxiety and cliff-edge dizziness, from mistaking heartbreak for cosmic crisis, and a place for suffering, each pain assigned a use in the manual. Many would give everything in the small hours for certainty. Here it takes three minutes.

State refusal's case fully too. Is the prescribed life still yours? If the machine recommends orchid-growing and following it truly makes you happy, where are you within that happiness? Might machine-given certainty remove precisely what makes happiness count? Think of the postcard: if a machine finishes and sends the sentence after ``when your summer holiday begins,'' what does Old Zhou actually receive?

Answer with reasons. There is no standard answer: here that is not politeness but the chapter's entire point. A world forbidding your own answer is precisely what all these philosophers warned against.

Finally, the postcard still lies on the table. Grandad is gone; Old Zhou has long been unreachable. Two and a half lines still stop at the summer holiday.

If you found it, what would you do? Put it back in the book, finish the sentence for him, or send it to the old address with a note explaining its history?

Your answer, and the way you justify it, constitute your response to this chapter's question.
